\chapter{Antecedentes}

La Universidad Nacional (UNA) es una institución pública de educación superior costarricense, comprometida con la excelencia académica y la pertinencia social de sus programas formativos. Desde su fundación en 1973, la UNA ha impulsado procesos de evaluación y mejora continua con el propósito de garantizar la calidad de la educación superior en el país \cite{UNAweb}.

En el ámbito de la acreditación, la Universidad participa activamente en los procesos regulados por el Sistema Nacional de Acreditación de la Educación Superior (SINAES), organismo que establece estándares de calidad a partir de dimensiones, componentes y criterios específicos \cite{SINAES}. Este marco exige a las carreras universitarias disponer de mecanismos formales para recopilar, organizar y reportar evidencias que sustenten el cumplimiento de dichos criterios.

Asimismo, la UNA administra de forma permanente recursos humanos y académicos a través de los denominados tiempos de jornada, los cuales representan unidades de planificación y presupuesto destinadas a la asignación de labores docentes y de proyectos institucionales. La gestión de este recurso resulta estratégica para garantizar la equidad en la distribución de la carga académica y la ejecución eficiente de iniciativas de investigación y extensión.

La revisión de documentos institucionales y normativas evidencia la necesidad de fortalecer los procesos administrativos y académicos mediante herramientas tecnológicas que permitan centralizar, automatizar y resguardar la información relacionada con acreditación y planificación de tiempos de jornada \cite{ReglamentoUNA, InformeSINAES}.
